\section{Related Work}
\label{sec:related}

\paragraph{Offline policy regularization.}
Offline RL methods control distribution shift either by constraining policy
improvement or by modifying the learned value function. Behavior constraints,
divergence penalties, and squared action anchoring keep the learned policy near
dataset support
\citep{fujimoto2019offpolicy,kumar2019bear,wu2019brac,fujimoto2021minimalist,tarasov2023revisiting}.
Other approaches regularize policy extraction through value-dependent
weighting or reshape the critic through pessimism and ensemble diversity
\citep{nair2020awac,kostrikov2022iql,kumar2020cql,an2021edac}. MPI addresses a
different control variable within TD3+BC: the movement of the actor branch
whose target copy enters critic learning.

\paragraph{Policy evaluation, target lag, and decoupled extraction.}
Approximate policy iteration conventionally alternates policy evaluation and
improvement, but their learned components can still share a target-policy
role. One-step offline RL instead freezes an estimate of the behavior policy's
value before a single constrained improvement
\citep{brandfonbrener2021onestep}. Backup-side constraints instead restrict
the actions selected inside the Bellman operator
\citep{kumar2019bear,ghasemipour2021emaq}, while decoupled extraction can train
a behavior model independently and use critic guidance only at inference
\citep{lin2026decoupling}. Polyak averaging delays a policy copy but retains
the same improvement branch for target construction and evaluation. MPI makes
a narrower distinction inside TD3+BC training: the first local refinement
supplies the target-policy branch and later refinements supply the evaluated
branch.

\paragraph{Trust regions, proximal methods, and policy geometry.}
Natural policy gradient and trust-region methods connect policy optimization
to Fisher information geometry \citep{amari2000methods,schulman2015trust}.
Optimal-transport geometry provides a distinct notion of policy displacement:
policy optimization has been formulated as a Wasserstein gradient flow, and
Wasserstein natural-gradient or trust-region methods have been developed for
continuous control \citep{zhang2018policy,moskovitz2020efficient,terpin2022trust,song2023metric}.
In offline RL, Wasserstein discrepancies have mainly been used as behavior
regularizers or distribution-matching objectives
\citep{wu2019brac,asadulaev2024rethinking}. Our use is different. The distance
defines an estimated-energy proximal operator; composing that operator with a
moving center defines a refinement path. We also distinguish the finite-step
variational statement from the small-step flow limit, rather than treating
large-budget end-to-end behavior as a direct numerical-integration test.

\paragraph{Expressive policy extraction and iterative refinement.}
Flow- and diffusion-based actors improve the expressivity of offline policy
extraction \citep{park2025flow,wang2023diffusion}. Their expressivity does not
by itself determine how aggressively a learned critic should move the policy,
and a single critic-guided extraction step can still be sensitive to
off-support value errors. MPI instead isolates sequential re-centering and the
policy branch used in critic feedback. The controlled experiments use
deterministic actors, for which the Wasserstein movement cost reduces to
squared action displacement.

\paragraph{Bootstrapping error and actor--critic coupling.}
Prior work has emphasized extrapolation and bootstrapping error as defining
failure modes of offline value learning \citep{fujimoto2019offpolicy,kumar2019bear,kumar2020cql}.
Our analysis focuses on an actor-side control variable that is usually
implicit: the displacement of the policy entering the TD target. The resulting
target-perturbation bound formalizes why target-policy displacement and final
deployment reach are distinct quantities, while the experiments measure their
realized counterparts under a jointly learned critic.
